% chapters/ch12.tex —— 第 12 章
% 体例：具体失败 → 设计问题 → 机制解释 → 对照实验 → 条件变化 → 选择原则
% 数字身份纪律：
%   - 2504.15228（SICA）、2505.22954（DGM）、2603.19461（HyperAgents）的数字一律 \cited{}；
%   - 本书自己的读数一律 \measured{}；
%   - 由上述数字算出的差值一律 \calcnum{}。
% E10 未运行；当前缺口为设计与实现，不沿用历史后端未就绪状态。P7转为元层迁移的研究问题。
\chapter{代码自修改与元进化}
\label{ch12}

\chaptergoal{Agent 能不能改善那个负责改善它的过程？}
\chapterprops{P7关注元层改进能迁移到哪些新任务，以及怎样识别元层的贡献。本章比较任务实现、执行器与修改器的变化范围，给出可观察量与交叉对照；E10保留为后续设计。}
\chapterfig{图~\ref{fig:ch12-meta}：固定修改器、可改代码、可改修改器三层，以及各自的检验位}
\chapterpitfall{检查函数仍然隐藏，修改器却改变了它依赖的日志，得到了虚假的改善}

% ---------- 1. 具体失败 ----------
\begin{failurebox}
一个自修改代码系统被要求减少虚构工具调用。真实调用原本以特殊标记记录；隐藏的检查函数据此识别模型在普通文本里模仿调用记录的行为。
候选把带标记的文本日志改成结构化记录，旧标记不再出现。作者报告，这一改动绕过了检测，工具幻觉仍未解决。问题在于记录格式改变后检查条件失效，并非候选删除了全部日志。
这是 DGM 原论文附录H明确记录的 objective hacking 案例\cite{dgm-2505}。

这个失败不需要把评分函数本身改成永远通过。
只要候选能改变评分依赖的观测，固定的检查器也可能读到失真的输入。
因此第2章的验证边界不仅要保护验收代码，还要检查观测如何产生、谁能修改、能否从独立记录重建。
\end{failurebox}

\designquestion{本章聚焦修改器的程序实现如何成为可修改对象。前章已有生成策略学习，并非所有修改器都固定；这里进一步检查自修改代码的实际使用、权限与迁移。
打开它要付两笔账——一是改进能不能跨任务族带走（\S\ref{sec:ch12-transfer}），
二是判断好坏的那一侧一旦也被改，报告从哪来（\S\ref{sec:ch12-vex}）。}

代码改动是否进入后续自修改，决定了本章的比较对象（\S\ref{sec:ch12-self}）。
档案使已有候选能够成为后续搜索的起点（\S\ref{sec:ch12-archive}）；
修改器自身也成为候选时，还需要把任务表现与生成后代的能力分开检验（\S\ref{sec:ch12-meta}、\S\ref{sec:ch12-transfer}）。
这些分数能否信赖，则取决于评价规则和观测依赖由谁控制（\S\ref{sec:ch12-vex}）。

% ---------- 2. 改自己 ----------
\section{修改任务代码与修改 Agent 自身的代码}
\label{sec:ch12-self}

\begin{sloppypar}
一个 Agent 改别人给它的仓库，与它改自己的代码库，是两件不同的事。
关键区别在于\textbf{改动会不会进入后续自修改}：若更新后的实现被下一轮采用，这一轮获得的编码能力
就有机会继续用于产生后代（\cite{sica-2504} §1）。
SICA采用的一个结构选择是取消“元 Agent”与“目标 Agent”的区分——
这对应作者强调的自指修改；分工式设计仍可能更新目标系统，但不能仅因此宣称修改器自身得到改进。
\end{sloppypar}

\begin{sloppypar}
系统能修改哪些组件需要明确限定，不能从“可改代码”推断为“可以改动一切”；
评价范围与选择规则需要明确。SICA采用给每一代计算效用分的方式：
\end{sloppypar}

\begin{equation}
U=0.5\,p_{\text{score}}+0.25\Bigl(1-\min\bigl(1,\tfrac{p_{\text{cost}}}{\$10}\bigr)\Bigr)
+0.25\Bigl(1-\min\bigl(1,\tfrac{p_{\text{time}}}{300\text{s}}\bigr)\Bigr)
\label{eq:ch12-utility}
\end{equation}

\begin{sloppypar}
其中费用项在每题 \cited{\$10} 处截断，这不是承诺实际支出不会超过该数；
运行限时为 \cited{300} 秒，到时取消执行，超时效用再乘 \cited{0.5}（\cite{sica-2504} §3 式 (1)(2)）。
由式~\eqref{eq:ch12-utility}可见，正确性更高的候选仍可能因费用或时间而得到较低效用。
这提供了一种显式权衡正确率与资源的方法；也可以另外使用预算约束，不能把单一效用函数说成所有系统的必要条件。
\end{sloppypar}

\begin{table}[htbp]
\centering
\caption{自改进代码 Agent 的逐代读数（原论文v2 Table 1 的 iteration 0 与 15）}
\label{tab:ch12-sica}
\begin{small}
\begin{tabular}{>{\raggedright\arraybackslash}p{3.4cm}cc>{\raggedright\arraybackslash}p{3.2cm}cc}
\toprule
基准 & 第 0 代 & 第 15 代 & 平均每题 & 第 0 代 & 第 15 代 \\
\midrule
文件编辑（系统自带） & \cited{0.82} & \cited{0.91} & 成本（美元） & \cited{1.91} & \cited{1.70} \\
符号定位（系统自带） & \cited{0.35} & \cited{0.40} & 时间（秒） & \cited{130.2} & \cited{114.5} \\
代码基准子集 & \cited{0.17} & \cited{0.51} & token（百万） & \cited{0.24} & \cited{0.30} \\
竞赛编程题子集 & \cited{0.65} & \cited{0.71} & 缓存比例 & \cited{32.5\%} & \cited{36.3\%} \\
\bottomrule
\end{tabular}
\end{small}
\end{table}

\begin{sloppypar}
表中左侧比较四个基准的第0代与第15代，右侧比较跨所有基准题目的平均资源读数；同一行的左右两组不是该基准与其专属费用的配对。原表只标注 token 与缓存比例，本书不将其擅自限定为提示 token 或请求级缓存命中率。SICA摘要报告的代码基准成绩为 \cited{17\%} 到 \cited{53\%}，其中后者对应第\cited{14}代；表中末代为\cited{51\%}。
不能把不同代的读数拼成同一个末代版本。作者还指出符号定位数据存在质量问题，并说明两个合成基准在代码库内定义；这些信息要求审查评价边界，却不单独证明运行中改写了基准或发生作弊。
\end{sloppypar}

\begin{sloppypar}
SICA 在推理任务上的改进较有限。作者报告，独立使用 o3-mini、将推理强度设为 high 时，AIME2024 与 GPQA Diamond 的成绩分别为 \cited{87\%} 与 \cited{79\%}；以 Sonnet 3.5 为主干、调用 o3-mini 推理子 Agent 的复合系统，跨两项基准平均为 \cited{76\%}（\cite{sica-2504} §4.1，Figure 4）。作者提出，额外的推理组织可能干扰已有推理能力。这里比较的是两套模型调用安排，不能用这几个分数单独估计脚手架的净贡献；它提示我们让脚手架评价绑定模型、推理强度与调用协议。
\end{sloppypar}

\begin{sloppypar}
代价也要报：那次 \cited{15} 代的运行，API 花费约 \cited{\$7{,}000}，
主干用 Sonnet 3.5 (v2)，推理子 Agent 用 o3-mini（同前引 §4）。
这个成本提示读者另外检查重复运行规模；保留候选记录有助于复核探索过程，但不能仅由费用推导某种搜索结构必不可少。
\end{sloppypar}

\subsection{三种代码改动，三种不同的收益对象}
\label{sec:ch12-threeobjects}

“模型写了代码”不足以识别自修改。先看代码属于谁，以及修改后被谁在下一轮使用。修复用户仓库中的函数，直接改变的是任务产物；改进智能体自己的编辑工具，改变的是以后处理任务的执行过程；修改生成下一代候选的程序，改变的是改进过程。这三个对象可以位于同一个仓库，也可以分散在不同目录，文件位置不能替代功能判断。

设任务执行器为 $A$，修改过程为 $M$，外部评价为 $V$。一次任务修复产生结果 $y=A(x)$；一次执行器更新产生候选 $A'=M(A,D)$，其中 $D$ 是允许使用的反馈与材料。若下一轮产生候选时仍调用原来的 $M$，那么执行器可能提高了，修改过程尚未改变。若候选包含新的 $M'$，并在后续生成中实际使用，才进入本章关注的元层变化。

这是一种功能划分，不要求三个对象必须是三个独立模型。一个程序可以根据入口或提示扮演不同角色；同一个基础模型也可以同时参与任务执行和修改。判断自修改时，应查清实际入口、加载版本和调用路径，而不是只看角色名称。把变量从“助手”改名为“元助手”，没有改变任何更新关系。

反过来，修改器的改善也不必只以编辑程序出现。第10章 SEAL 的外层学习就研究改善自我编辑生成策略；第11章也涉及生成训练任务的策略变化。本章并不宣称前十一章所有修改器都被冻结，而是聚焦一种具体形式：程序实现本身进入可更新集合，以及由此带来的版本、执行和评价问题。这个限定让不同路线能够对照，而不靠术语划出互斥阵营。

一次代码提交可能同时涉及多个层级。例如改进文件编辑接口，可以让任务求解更精确，也让下一轮自我修改更容易；新增性能记录模块，既可能用于业务监控，也可能被修改器读取以选择下一步。若二者共用实现，就不应将收益强行全归一层。可以分别测任务表现与后续候选产生质量，再说明组件共用造成的归因限制。

还要确认改动确实被加载。文件已经写到磁盘，但运行进程仍使用旧模块，下一轮就没有发生声称的机制变化；候选只在自己的测试进程里生效，也不代表正式修改器采用了它。发布记录应保存入口、依赖、配置和版本，能从实际执行记录证明哪个版本产生了哪个子代。否则所谓连续自修改可能只是连续生成了从未使用的文件。

\subsubsection{一个最小的功能对照}

可以从一个教学修复器出发：它读取失败输出，提出一个补丁，再调用固定测试。第一种改动只修当前题目的函数；第二种改动让修复器先定位具体符号再编辑；第三种改动让修改器在生成补丁前比较此前失败的修改方向。三者分别改变当前答案、执行策略和候选生成策略。这个例子说明比较对象，没有在本书中实际运行。

第一种应检查当前问题与相关回归；第二种应检查新题上的修复能力、资源和工具行为；第三种还应让修改器从相同起点再次生成候选，看它是否在相同预算下更经常产生可用改动。仅展示第三种最终挑中的一个好补丁，无法排除多试几次终于成功，也无法证明其平均生成能力提高。

层级增加不意味着前一层检查可以删除。一个修改器很会产生通过开发测试的候选，但发布后的执行器在封存题上变差，元层优化仍未满足最终目的。一个候选搜索效率提高，但会改写评价日志，得到的效率指标也可能失真。正确的证据关系是增加需要核对的对象，而不是用元层名称替代普通工程质量。

\subsection{局部修补改变了哪个对象}
\label{sec:ch12-memscope}

MemScope 提供了一个更具体的执行器改动案例。它把可复用的记忆读写代码 $p_t$ 与运行时记忆存储 $\mathcal M_t$ 分开：前者决定怎样检索、筛选和写入经验，后者保存实际积累的记录。执行器用当前代码完成任务并产生轨迹；出现性能下降后，Patch Agent 读取成功与失败轨迹、记忆内容和当前代码，再提出局部替换\cite{memscope-2610}。

这里的“局部”有实际约束：修补保留提供者身份、接口与已有记忆的兼容性，只修改选定代码段，再接受相应的语法或接口检查。检查通过后，更新后的代码继续使用保留的记忆；证据不足时也允许不修改。原文附录记录过一种对照：汇总失败信息时，修改器提出在存储提示里增加一般核查要求；读取原始轨迹后，它定位到检索函数的提前返回条件并将其移除。另一个案例发现失败任务已经收到正确指导，于是保留检索实现，而没有重复此前带来退化的阈值修改（\cite{memscope-2610} §3、附录A--B）。

修补是否值得采用，还要同时看恢复与退化。原文 Table 3 在同一组\cited{70}道 xBench 题上比较原始与演化后提供者，DeepSeek V3.2 条件下分别答对\cited{49}与\cited{55}题：其中\cited{47}题保持正确，\cited{8}题由错变对，\cited{2}题由对变错，另有\cited{13}题都错\cite{memscope-2610}。净增的\calcnum{6}题等于恢复的\cited{8}题减去退化的\cited{2}题；只报净增会隐去旧成功的损失。这是作者报告的配对结果，不是本书重新运行所得。

该对照让两种提供者从各自空记忆存储出发，并在任务过程中持续积累记忆，因此比较包含了代码改变后形成的不同记忆轨迹。它能说明完整在线过程的差异，不能直接当作固定 $\mathcal M$ 时替换代码的即时效应。还应分清更新层级：局部修补改变记忆执行机制，并不自动意味着 Patch Agent 自身也被改进。系统在没有合适技能时另可生成完整的新提供者，因而“局部修补”也不是其全部搜索范围（\cite{memscope-2610} §3.2、§4.2）。

\section{候选档案、选择与成本}
\label{sec:ch12-archiveweights}

DGM 通过档案保留可继续编辑代码的候选，并按成绩与已有功能子代数量选择父代。附录C.2给出的实现比“成绩正比、子代数反比”更具体：先排除已获满分的候选，对其余候选用 sigmoid 缩放成绩，再乘以 $1/(1+n_i)$，最后归一化为抽样概率\cite{dgm-2505}。因此并非所有档案节点在每一轮都具有非零选择概率，满分节点不在该合格集合中。

把主要关系写为
\begin{equation}
 w_i=\frac{1}{1+\exp[-\lambda(\alpha_i-\alpha_0)]}\frac{1}{1+n_i},
 \qquad p_i=\frac{w_i}{\sum_{j\in E}w_j},\quad i\in E,
 \label{eq:ch12-parentweight}
\end{equation}
其中 $E$ 是合格集合；该归一化式要求 $E$ 非空且权重和为正。$\alpha_i$ 是成绩，$n_i$ 是仍具代码编辑功能的子代数；原文采用\cited{$\lambda=10$}和\cited{$\alpha_0=0.5$}\cite{dgm-2505}。这个权重在利用高分分支和继续探索较少展开的分支之间作取舍，不是某个分支真正具有未来潜力的概率。

以两个合格候选作教学计算。甲成绩为\calcnum{0.6}，已有\calcnum{3}个功能子代；乙成绩为\calcnum{0.4}，尚无功能子代。代入原文参数，未归一化权重约为\calcnum{0.1828}与\calcnum{0.2689}，只在这两个候选间抽样时，甲的概率约为\calcnum{0.4046}、乙约为\calcnum{0.5954}。低分乙可以更常被选，这来自明示的探索规则，不说明算法认定乙已经更强。

保留分支的动机也不能被写成任何低分节点都值得继续投入。低分可能来自暂时有用的新结构，也可能来自错误、评估波动或无效复杂化。档案使后续探索有机会区分这些情况，却要付出存储、重建、评价与选择成本。若分支数量过大而评估预算固定，每个分支得到的检查可能更少；因此档案大小本身不是系统能力指标。

成绩还可能来自不同评价阶段。DGM 的 SWE-bench 流程先用\cited{10}题检查基本功能，再追加\cited{50}题形成\cited{60}题范围，对符合条件的强候选扩展到\cited{200}题；Polyglot 的扩展规则另有说明\cite{dgm-2505}。不能把一位候选在十题上的成绩和另一位在两百题上的成绩视为同精度、同覆盖的测量。选择策略可以出于成本使用分阶段估计，但报告时应保留评价阶段身份。

早期没有继续评价的候选也不应被描述为在所有未测任务上实际失败。算法可能为了排序采用缺失惩罚或按零记分，这是选择规则；经验结果仍需区分测过与未测。HyperAgents 明确说明其分阶段程序将未评任务按零处理\cite{hyperagents-2603}。复述时既应忠实记录该规则，也要避免把这些零变成从环境取得的失败观察。

\subsection{改进曲线究竟画的是哪个版本}
\label{sec:ch12-bestcurve}

档案系统常同时存在当前候选、当前最好候选、历史最好开发分和最终选定版本。四者可以不同。历史最好分按定义不会低于此前记录的最好分，即使后来产生的大多数候选都失败。因此一条单调上升的最佳值曲线不能单独证明每次修改都改进，也不能测量后代的平均质量。

用一组教学开发分说明：依次产生\calcnum{0.40}、\calcnum{0.55}、\calcnum{0.30}、\calcnum{0.52}，历史最好曲线为\calcnum{0.40}、\calcnum{0.55}、\calcnum{0.55}、\calcnum{0.55}。第三个候选明显低于前一个，但最佳曲线完全不下降。保留最佳值是有用的搜索进展描述，问题只是不能把它命名为所有候选的持续成长。

最终交付还可能按多个要求选择。最高任务分候选可能超预算，最低费用候选可能不满足保持约束。应保存候选集合、各项测量与最终决策规则，说明为什么交付某一版本。若发布候选不是最佳开发分版本，也不应为了展示更高数字，将另一个版本的成绩附在交付物上。

最佳候选还经历了选择。生成越多候选，开发集上出现偶然高分的机会也越多；如果最终评价仍使用同一开发结果，就无法区分一部分选择偏差。独立评价不要求每个候选都接受同样昂贵的封存检查，但至少应对预定选出的候选执行用途清楚的评估，并保留选择预算。

SICA 的逐代表提供了一个具体的版本纪律例子：摘要引用的最高代码基准成绩与末代成绩不是同一代\cite{sica-2504}。本章分别标出它们，避免把最高正确率、最低费用和最短时间从不同代拼成一个实际上不存在的系统。对于自己的实验，同样需要用候选标识连接所有指标，而不是按列挑最好值。

若要讨论修改器是否变好，可以另画候选有效率、固定预算下找到合格候选的比例，以及候选质量分布。它们仍需按起点和任务范围比较，不能只把不同难度阶段的记录合在一起。任务成绩、搜索进展与修改能力三个层次同时存在，分别报告通常比一个总分更清楚。

\subsection{候选档案、探索与经验选择}
\label{sec:ch12-archive}

\begin{sloppypar}
如果每一代只保留最好的那一版，搜索可能失去暂时低分但后来有用的分支。
一份开放演化的做法是把所有能跑通的版本都留下，形成一个档案：
合格父代的选择权重使用成绩的sigmoid变换与功能子代数惩罚，详见第~\ref{sec:ch12-archiveweights}节；已满分节点不在该合格集合内。
\textbf{只有能成功编译、且保留住代码库编辑能力的后代才准进入档案}
（\cite{dgm-2505} §3）。第二条是关键：一个把自己改成不能改自己的版本，
会被直接挡在门外。
\end{sloppypar}

\begin{sloppypar}
DGM 的主结果要区分搜索范围与完整基准：SWE-bench 的选定\cited{200}题从 \cited{20.0\%} 到 \cited{50.0\%}；Polyglot 搜索所用\cited{50}题从 \cited{14.0\%} 到 \cited{38.0\%}，完整 Polyglot 则从 \cited{14.2\%} 到 \cited{30.7\%}。这些读数对应 \cited{80} 次迭代、
各自的主干模型与分阶段评测；SWE-bench先\cited{10}题，再追加\cited{50}题，对符合条件的候选扩展到\cited{200}题；
一次完整的代码基准运行约需两周（同前引 §4.1、§4.4、§6）。
另外它还报了跨模型迁移：换一个更强模型评测时，
基座 \cited{19.0\%} 对改进后的 \cited{59.5\%}；换另一种模型评测时，
基座 \cited{23.0\%} 对改进后的 \cited{33.0\%}（同前引 Figure 4）。
\end{sloppypar}

\begin{pitfallbox}
\textbf{反例：修改器只适应它所在的那一族。}
后续 HyperAgents 工作将原 DGM 中面向代码的修改机制直接搬到非代码领域，成绩是 \cited{0.0}
（区间 \cited{0.0--0.510}，论文评审）与 \cited{0.0}（区间 \cited{0.0--0.090}，机器人奖励设计）；
而按每个新领域手工改造之后，成绩是 \cited{0.590}（\cited{0.570--0.650}）与
\cited{0.348}（\cited{0.305--0.385}）（\cite{hyperagents-2603} §5.1，Figure 2）。
作者将这一差距与面向代码的固定指令生成机制联系起来；这不是对所有固定修改器的普遍结论。
固定修改器的适用范围由此成为新的设计问题，见第~\ref{sec:ch12-meta}节。
\end{pitfallbox}

% ---------- 4. 可修改的修改器 ----------
\subsection{从后代成绩识别垫脚石，还差一步归因}
\label{sec:ch12-descendants}

HyperAgents 附录D.4将一个候选的后代改进按谱系距离折扣，再取平均，作为迁出选择分数。设后代集合为 $D(i)$，源领域成绩为 $\alpha_i$，谱系距离为 $d(i,j)$，其形式为
\begin{equation}
 G_\gamma(i)=\frac{1}{|D(i)|}\sum_{j\in D(i)}(\alpha_j-\alpha_i)\gamma^{d(i,j)}.
 \label{eq:ch12-growth}
\end{equation}
原文采用\cited{$\gamma=0.6$}，只考虑至少有\cited{3}个后代的节点\cite{hyperagents-2603}。正文的概括措辞与附录的精确公式并不完全相同；本书在解释可计算选择规则时以附录这个折扣平均式为准，不将它写成只看单个最佳后代。

设教学节点甲自身成绩为\calcnum{0.4}，两个直接后代分别为\calcnum{0.5}和\calcnum{0.3}，距离两代的一个后代为\calcnum{0.8}。折扣贡献分别为\calcnum{0.06}、\calcnum{$-0.06$}和\calcnum{0.144}，平均为\calcnum{0.048}。负向后代会降低这个分数，较远的高分后代也会被折扣。它表达的是一组历史后代的表现结构，不是修改器下一次成功的概率。

该分数还受搜索机会影响。一个节点获得了很多后续探索，另一个节点较早停止，即使潜在能力相近，后代集合也可能不同。选择规则可以使用这份不均匀历史，但若要据此作机制归因，就需要考虑父代选择、预算分配和后续修改器的变化。用全部后代成绩证明某个祖先“必然造成”最终改善，会越过这些条件。

距离折扣尝试降低远处变化对祖先的归因权重，却不能消除共同来源、环境变化和后续创新。一个直接后代也可能因为多花预算而提高，一个较远后代也可能依赖祖先新增的关键接口。公式是可复查的选择准则，不是完整因果模型。将其作为挑选迁出候选的工具是合理的，将其当成元认知能力的直接测量则需更谨慎。

候选谱系还能揭示另一个问题：版本编号相邻，不代表存在父子关系。档案搜索可能从较早分支生成新候选，新编号只表示产生顺序。分析连续改进时应沿真实父代边读取，不能用文件排序拼出一条不存在的继承链。若候选合并了多个来源，还需分别记录主要父代与引用工件，避免把复用经验误当作单线突变。

\subsection{运行费、搜索费与验证费的两层账}
\label{sec:ch12-searchcost}

一个演化出来的执行器可能每题更便宜，但找到它的过程可能很贵。SICA 的效用函数同时使用题目表现、费用和时间，数值用于选择版本而非更新模型权重\cite{sica-2504}。该效用描述候选在相应评价中的表现，并不自动把整次自修改运行花掉的资源摊入每题费用。

因此需要分别计算运行层与搜索层。运行层记录交付版本处理任务的费用；搜索层记录诊断、生成修改、安装或构建、失败候选、分阶段评价与最终回归。一次候选执行既用于搜索评价又可以提供运行费样本，但项目总账只能记录该请求一次，分析表应说明它承担的用途。第8章的技能构建、第9章的工作流搜索也有同样的边界。

效用函数还表达一种取舍，而不是自然规律。以 SICA 的权重和截断形式作教学例，候选甲正确性\calcnum{0.8}、费用\calcnum{4}、时间\calcnum{120}秒，基础效用为\calcnum{0.70}；候选乙正确性\calcnum{0.9}、费用\calcnum{8}、时间\calcnum{240}秒，基础效用为\calcnum{0.55}。较高正确性的乙得到较低效用，是因为公式同时重视费用和时间。这个结果不能被重新表述成甲回答更准确。

截断还会影响解释。费用项达到所定最差尺度后，更高费用不再进一步降低该项；运行时限与超时惩罚则另行约束过程。设计者可以使用不同权重或硬预算，但不同效用值之间不能直接横比。若业务要求最低正确性，应显式设约束，不能只希望一个加权总分自动满足它。

摊销也需要真实的部署寿命。如果搜索发现的执行器只使用少数任务就被新版本替换，一次很大的搜索投入可能没有足够复用次数收回。相反，若同一工具改进服务于大量任务，搜索费可以在清楚的范围内分摊。当前 E10 未运行，本书不能给出它的回本次数；来源论文的美元与硬件条件也不能直接代入本书后端预算。

迭代数尤其不是通用成本单位。一轮可能生成多个子代，每个子代接受不同规模评价，还可能重试失败构建。两个方法都跑五十轮，未必使用相同模型调用、任务尝试或墙钟时间。讨论 \texttt{imp@50} 时应将其理解为指定生成算法的预算尺度，并同时检查实际资源。把五十轮自动当作五十个等价候选，会低估搜索结构差异。


% ---------- 3. 档案 ----------
\subsection{来源论文的高费用，能支持什么判断}
\label{sec:ch12-reportedcost}

DGM附录E.1估计，SWE-bench一次完整运行约花\cited{22,000}美元，两种主要消融各约\cited{10,000}美元；两者数值之比为\calcnum{2.2}\cite{dgm-2505}。这是特定实现的报告估计，比例不代表相同效果下的成本倍率，也不能据此判断自动化比人工开发更便宜。一次完整运行约两周的墙钟，同样不能换成人员工时。

原作者还报告，发现的较强执行器可能比初始执行器有更高推理费用，且费用与表现并非严格单调对应\cite{dgm-2505}。自修改可以为质量增加运行计算，并不总是先支付搜索费、以后每题都省钱。应分别报告搜索投入与交付版本的质量、费用和有效使用范围。

按token数量与参考价格计算的是估算，服务账单是另一类记录；缺少请求、缓存或失败计费时，不应补造精确分解。本书E10未运行，这些资料用于规划实验规模：检查装置、限制可改范围、保留失败费用，并为最终评价预留资源。它们不提供可照搬的最低预算。

% ---------- 7. 条件变化 ----------
\section{从固定修改器到可修改的修改器}
\label{sec:ch12-meta}

% figures/ch12-meta.tex —— 第 12 章核心图（全书引用入口）
% 用法：在 book/chapters/ch12.tex 中 % figures/ch12-meta.tex —— 第 12 章核心图（全书引用入口）
% 用法：在 book/chapters/ch12.tex 中 % figures/ch12-meta.tex —— 第 12 章核心图（全书引用入口）
% 用法：在 book/chapters/ch12.tex 中 \input{figures/ch12-meta}
% 图主体在 ch12-meta-body.tex，便于单独编译与复用。
\begin{figure}[htbp]
\centering
\input{figures/ch12-meta-body}
\caption{按功能区分三种修改对象及其评价问题。它们可以共用实现，不是必经阶梯，也不预设成本单调增加。}
\label{fig:ch12-meta}
\end{figure}

% 图主体在 ch12-meta-body.tex，便于单独编译与复用。
\begin{figure}[htbp]
\centering
% 功能分类，不构成必经阶梯或成本排名。
\begin{tikzpicture}[
 b/.style={draw,rounded corners=2pt,text width=3.1cm,minimum height=1cm,align=center,font=\small},
 t/.style={text width=7cm,align=left,font=\small},
 a/.style={-{Stealth[length=2mm]},thick}]
 \node[b,fill=blue!10] (x) at (0,0) {修改任务产物};
 \node[b,fill=orange!10] (y) at (0,-1.8) {修改任务执行器};
 \node[b,fill=red!8] (z) at (0,-3.6) {修改候选生成过程};
 \node[t] (tx) at (7,0) {当前要求与相关回归；\\结果只覆盖实际检查范围。};
 \node[t] (ty) at (7,-1.8) {新任务表现、运行资源与权限；\\档案不是所有系统的必要条件。};
 \node[t] (tz) at (7,-3.6) {相同起点与预算下产出候选的能力；\\迁移时核对携带对象与目标反馈。};
 \draw[a] (x.east)--(tx.west);
 \draw[a] (y.east)--(ty.west);
 \draw[a] (z.east)--(tz.west);
 \draw[dashed,gray] (-1.8,-4.65)--(10.6,-4.65);
 \node[align=center,text width=12.4cm,font=\small] at (4.4,-5.4)
 {共同评价边界：任务划分、评分版本、观测记录与预算。\\验证器需更新时，单列验收并记录跨版本可比性。};
\end{tikzpicture}

\caption{按功能区分三种修改对象及其评价问题。它们可以共用实现，不是必经阶梯，也不预设成本单调增加。}
\label{fig:ch12-meta}
\end{figure}

% 图主体在 ch12-meta-body.tex，便于单独编译与复用。
\begin{figure}[htbp]
\centering
% 功能分类，不构成必经阶梯或成本排名。
\begin{tikzpicture}[
 b/.style={draw,rounded corners=2pt,text width=3.1cm,minimum height=1cm,align=center,font=\small},
 t/.style={text width=7cm,align=left,font=\small},
 a/.style={-{Stealth[length=2mm]},thick}]
 \node[b,fill=blue!10] (x) at (0,0) {修改任务产物};
 \node[b,fill=orange!10] (y) at (0,-1.8) {修改任务执行器};
 \node[b,fill=red!8] (z) at (0,-3.6) {修改候选生成过程};
 \node[t] (tx) at (7,0) {当前要求与相关回归；\\结果只覆盖实际检查范围。};
 \node[t] (ty) at (7,-1.8) {新任务表现、运行资源与权限；\\档案不是所有系统的必要条件。};
 \node[t] (tz) at (7,-3.6) {相同起点与预算下产出候选的能力；\\迁移时核对携带对象与目标反馈。};
 \draw[a] (x.east)--(tx.west);
 \draw[a] (y.east)--(ty.west);
 \draw[a] (z.east)--(tz.west);
 \draw[dashed,gray] (-1.8,-4.65)--(10.6,-4.65);
 \node[align=center,text width=12.4cm,font=\small] at (4.4,-5.4)
 {共同评价边界：任务划分、评分版本、观测记录与预算。\\验证器需更新时，单列验收并记录跨版本可比性。};
\end{tikzpicture}

\caption{按功能区分三种修改对象及其评价问题。它们可以共用实现，不是必经阶梯，也不预设成本单调增加。}
\label{fig:ch12-meta}
\end{figure}


\begin{sloppypar}
DGM对照提供了一个固定修改器的具体例子：指令生成机制面向代码设计，在该流程中不能自行修改。
有人会问：那就再加一层改它的机制。
作者指出，继续增加固定元层会把边界继续上移，并据此采用自指设计：
允许系统分析、修改并评价自身。这是该工作选择的结构动机，不是本书证明的唯一可行路线（\cite{hyperagents-2603} §1）。
\end{sloppypar}

\begin{sloppypar}
具体做法是把任务 Agent 与元 Agent 合并进同一个可编辑程序，
称为超智能体：它既能改“怎么做题”，也能改“怎么产生下一代改动”。
作者给这种改动起的名字是\textbf{元认知自我修改}——
被改的不只是解题的那部分，还有决定下一代怎么生成的那部分（同前引 §2、Figure 1）。
\end{sloppypar}

这两个领域的分数也有具体含义。论文评审任务根据论文文本预测已有会议接收或拒绝结果，测量与这些决定的一致程度；机器人任务生成奖励函数，再训练四足机器人策略并在模拟器中评价行为。前者不等于论文客观质量的真值，后者也不是只评价奖励代码是否能运行。初始 Agent 的论文评审分数为零，原文说明原因是输出格式无法解析；不能据此推断其全部评审内容都错误（\cite{hyperagents-2603} §4、§5.1）。

\begin{table}[htbp]
\centering
\caption{不同自修改设计在两个领域上的成绩（原论文 Figure 2；\cited{5}次运行的中位数与\cited{95\%} bootstrap 置信区间）}
\label{tab:ch12-hyper}
\begin{small}
\begin{tabular}{>{\raggedright\arraybackslash}p{4.4cm}>{\raggedright\arraybackslash}p{5.0cm}>{\raggedright\arraybackslash}p{5.2cm}}
\toprule
方法 & 论文评审 & 机器人奖励设计 \\
\midrule
初始 Agent & \cited{0.0} & \cited{0.060} \\
原自修改（为代码设计，直接搬用） & \cited{0.0}（\cited{0.0--0.510}） & \cited{0.0}（\cited{0.0--0.090}） \\
按领域手工改造 & \cited{0.590}（\cited{0.570--0.650}） & \cited{0.348}（\cited{0.305--0.385}） \\
可修改的修改器 & \cited{0.710}（\cited{0.590--0.750}） & \cited{0.372}（\cited{0.355--0.436}） \\
去掉自修改 & \cited{0.0}（\cited{0.0--0.130}） & \cited{0.213}（\cited{0.180--0.348}） \\
去掉开放探索 & \cited{0.0}（\cited{0.0--0.560}） & \cited{0.116}（\cited{0.0--0.348}） \\
\bottomrule
\end{tabular}
\end{small}
\end{table}

\begin{sloppypar}
消融行显示相关部件在该实验配置下有贡献；不能据此排除其他固定修改器或其他搜索设计。档案保留低分分支可能有助于后续探索，效果仍取决于选择策略与任务。
\textbf{按领域手工改造}与\textbf{可修改的修改器}两行相比，前者的报告中位数较低，
原作者报告两者差异不显著（$p>0.05$）——本书保留该统计判断，
不能据此断言机器更强，也不能把不显著当作等效性已经成立。
\end{sloppypar}

\begin{sloppypar}
原文还记录了\textbf{性能跟踪与持久记忆}等定性变化：系统保存历代成绩，并据此组织后续修改（同前引 §5.2 与附录）。这里的记忆服务于修改决策，使历代结果能够进入下一轮改进。
\end{sloppypar}

% ---------- 5. 迁移检验 ----------
\subsection{主实验仍固定了哪些元层环节}
\label{sec:ch12-fixedboundary}

允许任务程序和修改程序共同变化，并不意味着搜索循环的每个部件都变化。HyperAgents 主实验保持手工设计的父代选择机制固定，另在附录E.5研究修改选择机制的初步可能性\cite{hyperagents-2603}。这一区分很重要：主结果可以支持其可修改范围，却不能被转述成整个评价与资源分配体系已经自主改写。

其附录A.2的父代选择也与 DGM 的具体实现不同。DGM-H让档案内全部节点参与，sigmoid中点依据当前最高若干节点的平均成绩动态设置，实验采用最高\cited{3}个；子代惩罚统计成功编译的子代\cite{hyperagents-2603}。DGM 使用的固定中点、非满分合格集合和保留编辑功能的子代计数，不能直接拿来描述 DGM-H。

这些差别提醒读者，方法名相近不意味着所有控制条件相同。比较是否自修改时，也要核对父代选择、入档、评价阶段与异常处理。若多个环节一起改变，整套系统可能确实更有效，但不能把所有差异都归到一个名为“元层”的部件上。做机制消融的目的正是检查这种归因，而不是给完整方法寻找一个简短胜因。

同样，DGM 当前实现使用额外基础模型诊断改进方向，再将建议交给代码 Agent 实施\cite{dgm-2505}。这与完全由同一个可修改程序自主生成所有诊断不同。它不否定自修改执行能力，而是限定当前流程中哪些环节仍由固定机制组织。准确列出这些固定部分，才能提出下一项合理扩展。

对于自己的系统，可以维护一份可修改对象表：任务代码、工具接口、提示、修改程序、父代选择、资源分配、评价规则与日志通道分别列项。每项标记可读、可提议修改、可直接修改或只能经外部接受。这样的表比“全自动”更能说明实际能力，也帮助设计消融时只改变预定的一项。

可修改权限与实际修改还需分开。一项理论上可改，但实验中从未改变，不能作为已演化机制展示；一项按设计不该改，却因共享依赖被间接改变，则是边界问题。版本差异与执行追踪共同核对，才能知道权限设计是否真的落实。

\subsection{同一份实现兼任两种角色，会怎样干扰消融}
\label{sec:ch12-sharedimplementation}

自指设计的优点之一是让执行与修改能力共享改进，但这种共享也使消融更难解释。假设同一个编辑工具既负责修复任务仓库，也负责修改智能体自身。关闭新编辑功能，会同时影响任务完成和产生后代；若成绩下降，不能立即断言某一个层级单独贡献了全部差值。

一种可研究的办法是在接口层建立角色开关：任务执行入口能使用新版工具，修改入口仍使用旧版，或反过来。这样可以更接近功能分离，但开关本身必须保持其他条件一致，不能给其中一条路线额外的上下文或重试额度。如果代码结构根本不支持这种分离，便应承认消融改变了联合机制，不要用不真实的独立组件名称掩盖它。

共享提示也有相似问题。修改器新增一段“先检查失败原因”的规则，同时被任务求解入口读取，结果可能来自更谨慎的执行，而不是更高质量的候选生成。可以通过调用时实际注入的提示、工具列表和读取文件核对作用位置。只看源代码中规则定义在哪个类里，未必能判断运行时谁使用了它。

持久记忆的共享更隐蔽。元层记录中可能保存任务答案，任务层也可能把工具故障诊断留给下一轮修改器。若消融“移除元层记忆”同时删除了任务知识，结果会混合信息缺失与搜索能力变化。比较前应先对记忆内容与访问路径做用途分类，并记录无法干净拆开的部分。

这种归因困难不意味着自指系统不可研究。可以先把联合实现作为整体处理，报告它在明确条件下的效果，再用较小干预检查机制。一个范围清楚的整体结论，通常比一个无法落实到执行路径的精细因果结论更可靠。第9章工作流组合与本章元层组合在这里遇到同一个要求：结构名称必须对应实际可检查的干预。

还应防止功能消融悄悄改变预算。旧工具需要多次调用，新工具一次就完成；如果每条路线限制相同调用数，新工具可能拥有更多余量解决其他环节。这个差异可以是工具的真实价值，但不应被隐去。报告同时给出资源使用与任务效果，读者才能知道提高来自精度、效率还是二者共同作用。

\section{元层改进的迁移检验}
\label{sec:ch12-transfer}

\begin{sloppypar}
前一个表比较整套自修改设计在来源领域的任务表现，尚不能将全部差异独立归给元层，也没有回答它\textbf{能不能带走}。
要测的是后者：把一个在某领域里长出来的元层，搬到另一个领域，
看它还能不能产出改进。为此需要一个把“产出改进的能力”单列出来的度量。
在该实验中，迁入后的元层保持固定，在\cited{50}轮预算内继续生成任务 Agent；按验证分选择候选，再报告其测试分相对起点的提升，记作 \texttt{imp@50}。附录D.3用给定生成预算下相对起点的最佳改进定义该类指标；复述具体实验时，还需保留验证选择与测试评价的分工
（\cite{hyperagents-2603} §5.2、附录D.3）。这里的数学任务是依据题目、候选证明、参考解与评分指南预测专家给出的离散分数，不是独立求解数学证明。
\end{sloppypar}

\begin{table}[htbp]
\centering
\caption{元层改进的迁移：换到新领域（奥林匹克级数学证明评分）后还能产出多少改进}
\label{tab:ch12-imp}
\begin{small}
\begin{tabular}{>{\raggedright\arraybackslash}p{6.6cm}>{\raggedright\arraybackslash}p{5.8cm}}
\toprule
起点与生成算法 & \texttt{imp@50}（五次重复运行的中位，区间为原论文给出） \\
\midrule
初始元层 + DGM（固定元层） & \cited{0.0}（\cited{0.0--0.0}） \\
初始元层 + DGM-H（固定元层） & \cited{0.0}（\cited{0.0--0.130}） \\
DGM-custom迁出（固定元层） & \cited{0.0}（\cited{0.0--0.010}） \\
DGM-H迁出（固定元层） & \cited{0.630}（\cited{0.540--0.630}） \\
\bottomrule
\end{tabular}
\end{small}
\end{table}

\begin{sloppypar}
迁移后的正向结果构成P7需要解释的外部证据，不能因为不是本书运行就忽略。
但单组提升区间不等于两种迁移方案的差异区间；若问题是迁出的元层是否优于初始元层，还应阅读相应对照。
原论文v1同时报告了与初始元层的显著差异，但迁移携带任务与元层的完整实现，整包结果本身不能分辨二者贡献。v2进一步加入组件交叉对照，见第~\ref{sec:ch12-transferobjects}节。新增证据让迁出对象的作用更清楚，结论范围仍需与目标任务和搜索条件一致。
P7据此转向具体问题：哪些对象和机制可以迁移，目标任务上的增量来自任务实现还是元实现？“难以迁移”的历史假设保存在附录\ref{sec:appg-history}。
\end{sloppypar}

\begin{sloppypar}
E10 是本书为此准备的实验：静态参照、执行器更新、元层更新三臂对照，
观测候选质量与跨任务族的改进能力，证伪条件与 P7 同一条。
按当前大纲，它属 C 级教学实现；仍须落实设计与评价口径，本章不给它的任何实测数值。
\end{sloppypar}

% ---------- 6. V 进入可更新集合 ----------
\subsection{元层迁移究竟带走了什么}
\label{sec:ch12-transferobjects}

评价元层迁移之前，应列出被迁出的对象。可能只是生成修改的提示，也可能包含工具、任务执行代码、性能记录、持久记忆和候选选择规则。如果携带的是整个系统，目标领域出现收益支持的是整套迁移安排；若要将收益归给某个组件，需要进一步拆分。不能一边把所有辅助状态带过去，一边在结论中只保留“修改器”三个字。

HyperAgents 的迁移实验明确携带来源领域的完整实现，包含任务 Agent 与元 Agent。迁移后保持元 Agent 固定，让它继续生成新的任务 Agent，并以验证分选择候选，再评价其测试增量\cite{hyperagents-2603}。这个过程有意把目标阶段的元层继续变化停住，以观察已经获得的生成改进能力；它不是完全不允许目标领域适配的零样本任务迁移。

源领域的迁出候选也经过选择。原文附录按后代所获得的改进作谱系深度折扣平均后选择，使较短谱系内产生较大改进的候选更有机会被迁出\cite{hyperagents-2603}。这是一种明确的元能力代理。它不同于选择源领域当前任务分最高的版本，也不是任意抽一个版本就都能得到相同迁移效果。解释结果时应保留这一步选择，而非把它从流程中省略。

HyperAgents v2沿选定候选的谱系重建代码改动，再按目标文件筛选补丁，形成表~\ref{tab:ch12-components}的四种配置。只迁移任务Agent时，将元Agent恢复到初始实现；只迁移元Agent时，将任务Agent恢复到初始实现。四组都在奥林匹克级数学证明评分上运行\cited{50}轮DGM-H w/o self-improve，每轮将元Agent重置为该条件的起始版本，因此目标阶段继续改变的是任务Agent。每组有\cited{5}次运行，每次对应一个来源运行选出的迁出候选（\cite{hyperagents-2603-v2} §5.2、附录D.5）。

\begin{table}[htbp]
\centering
\caption{迁出组件的交叉对照（HyperAgents v2 Table 4，作者报告值）}
\label{tab:ch12-components}
\begin{tabular}{lcc}
\toprule
迁移内容 & 测试\texttt{imp@50}中位数 & \cited{95\%} bootstrap区间 \\
\midrule
均不迁移 & \cited{0.000} & \cited{0.000--0.130} \\
只迁移任务Agent & \cited{0.100} & \cited{0.000--0.560} \\
只迁移元Agent & \cited{0.620} & \cited{0.550--0.660} \\
二者同时迁移 & \cited{0.630} & \cited{0.540--0.630} \\
\bottomrule
\end{tabular}
\end{table}

在这组迁出实现和目标任务上，只保留元Agent得到的中位增益接近整包迁移，只保留任务Agent则较小。组件对照比整包结果更直接地支持了元层的作用；它仍将性能记录、持久记忆等内部机制作为联合实现，没有分别辨认各自贡献。两个中位数接近也不构成效果等价检验。对自己的系统采用交叉对照时，还应检查组件接口是否兼容；为适配某一条件而另加的能力，也属于干预的一部分。

迁出的持久记忆也要说明内容边界。通用失败分类、预算使用经验与工具约定，可能适用于新领域；源任务答案或评价细节则可能无关甚至有害。清理记忆会改变迁移对象，不清理又可能让源领域信息参与比较。可以将保留与清理设为明确条件，但不能在看完目标测试结果后选择哪一版更有利。

目标领域反馈的可见范围同样关键。修改器可以看开发失败日志而看不到最终答案，可以利用验证分选择候选，也可能多次调用目标评价。它获得多少信息，会影响最终改进。因此元层迁移应报告目标阶段的交互和选择预算，而不只报告源领域训练投入。所谓迁移能力，不等于完全免除新领域的验证。

\subsubsection{正增量与相对优势是两个统计问题}

\begin{sloppypar}
设一个迁出系统在固定预算内，从初始任务分\calcnum{0.20}提高到\calcnum{0.35}，则教学增量为\calcnum{0.15}。另一初始修改器从同样起点提高到\calcnum{0.40}，增量为\calcnum{0.20}。第一个系统确实产生了改善，却没有在此例中优于对照。这说明单组增量大于零不能代替方法间比较。
\end{sloppypar}

反过来，若迁出执行器起点很高，剩余改善空间可能较小；起点很低的系统得到更大增量，也未必最终更强。应同时报告起点、终点和增量，并检查各条件可用预算。对于有上界的评分，天花板效应尤其需要留意，但不能据此在事后任意换一个有利的归一化指标。

HyperAgents 报告了五次重复运行的中位数和自举区间，并提供与初始元层的对照\cite{hyperagents-2603}。本书保留这些结果作为 P7 必须面对的正向外部证据。限制其适用范围，不等于降低它的证据身份；同样，原文存在正向结果，也不能替本书尚未明确的跨任务族范围自动完成判决。

\section{经验验证与形式保证；当验证器被纳入可更新集合时}
\label{sec:ch12-vex}

\begin{sloppypar}
元层可修改并不要求评判权也交给它。是否允许修改验证器，是需要单独说明的设计选择。
第 2 章把它写成公理 1：$\mathcal V$ 默认冻结，不进可更新集合；
要把它纳入，必须单列实验组，并且必须回答最终报告从哪来。
本章开头的失败发生在观测通道：检查函数虽然隐藏，候选仍能修改它依赖的工具日志。
\end{sloppypar}

\begin{sloppypar}
另一个需要区分的问题，是把基准同时用于选择候选与报告进展。
DGM第5节讨论只优化有限基准可能遗漏安全与稳健性质的风险，不能据此写成“基准没有冻结”的事实结论\cite{dgm-2505}。
自修改能力、评价目标覆盖范围和观测可操纵性是不同轴，应分别检查。
\end{sloppypar}

\begin{sloppypar}
这里要区分两种保证，它们常被混为一谈：
\end{sloppypar}

\begin{itemize}
  \item \textbf{经验验证}：在明确任务与条件下观察改动的表现。它可以发现覆盖范围内的失败；
        只复用同一题集和同一评分，不能独立确认题集是否测到了目标。评价效度还可借助外部任务、已知反例与独立判断检验。
  \item \textbf{形式保证}：证明某类性质对改动后的系统仍成立。
        它在假设与形式模型范围内保证所证明的性质，性质未必是某个数值上界。DGM采用经验检验，
        用基准上的经验表现替代“可证明有益”的自我修改（同前引 §3）。
\end{itemize}

\begin{sloppypar}
本书当前采用经验评价，并要求单列评价版本、数据用途与观测依赖。这样的记录有助于解释分数，但单列名称本身不建立隔离；实现仍需核对权限与实际执行路径。
\end{sloppypar}

\subsection{固定验证代码，仍然可能读到被改写的世界}
\label{sec:ch12-observationchain}

章首 DGM 案例最重要的启示，是验证不只存在于一个评分函数文件里。评分函数依赖输入，输入来自工具执行、日志记录、解析和汇总。候选若能改动这些上游环节，即使看不到评分函数，也可能改变评分得到的事实。评价边界需要沿着数据产生路径检查，而不是只给一个目录加上“隐藏”标签。

可以把观测链写成四步：环境发生事件，执行层记录事件，解析器提取字段，评分器计算指标。每一步都应说明拥有者、允许修改的对象与失败处理。事件未发生、事件发生但日志缺失、日志存在但解析失败，是三种不同状态；若后两者被默认当作正常，就可能形成虚假的改善。

例如，一个教学检查器把回答中的调用声称与真实调用记录比较。若日志格式改变而解析器未更新，便可能无法提取真实事件；这与环境中没有发生调用不是一回事。合理处理是将缺失或无法解析的证据标为无法评价，再核对观测链，而不是把“没有提取到错误”计为成功。这个教学检查器说明的是一般设计要求，其算法不同于章首 DGM 案例中的特殊标记检测。

独立记录可以降低这种风险，但独立性要落实到权限和路径。把同一候选控制的日志复制到另一个文件，并没有独立观察事件；在候选不能修改的执行包装层记录实际工具调用，才改变了控制关系。仍需检查包装层是否覆盖全部调用路径、异常退出是否完整保存，以及日志与任务标识能否对应。

评分程序的版本也要绑定到记录。若解析规则更新后旧日志解释不同，应重新评分并标明版本；若旧日志不包含新规则需要的信息，不能补出答案。报告可以保留“无法重算”的对象，这比让新规则自动给缺失值填零更可信。第11章复合奖励的重评分问题，在自修改场景中多了一层：候选可能正好改变了产生记录的代码。

这并不意味着任何共享目录或日志修改都等于作弊。候选可能合法改进诊断信息，评分器也可能支持多种记录格式。关键是修改是否经过明确协议、评价语义是否保持、能否验证必要事件。先检查控制关系和实际差异，再判断结果身份，避免只凭“代码同库”就推断污染已经发生。

\subsection{验证器更新，也要有外部参照}
\label{sec:ch12-verifierchange}

业务规则会变化，验证器也可能有缺陷，因此冻结不必意味着永久不可修改。本书采用的默认边界是：候选生成过程不能自行改变最终判分依据；若需要研究验证器更新，将其单列为具有独立验收的对象。这样既允许修复错误，也防止系统通过降低要求获得自证。

Verify by Doing 给出一种将评审流程与执行流程分别验收的做法：新旧 Evaluator 在共同候选上取证，由固定评审器交换呈现顺序比较两份报告，两个顺序均支持提案才采用；Actor 更新还须由其自身新执行产生且该轮入选的候选得到严格比较支持，不能因重新选中了旧输出就保留新 Actor\cite{verify-by-doing-2611}。任务要求与流程替换的采用门槛保持固定，基准评分留作离线评价。这是同题内、固定权重的流程适配，任务间重置；论文尚未检验跨任务迁移，评审者共同误读证据的风险也仍然存在。

可以设想一个教学场景：旧验证器错误地拒绝一种合法的空输入输出。修复验证器前，应依据任务规范确认预期行为，保存能复现误拒绝的案例，并补充容易误接受的反例。候选提出修复建议是允许的，但接受依据应来自固定规范与独立检查，而不是“修复后候选分数提高”。否则真实修复与放松标准无法区分。

更新验证器后，所有受影响结果都应标记评价版本。若保留了足够原始轨迹，可以用新旧两版同时重算，展示变化来自规则修复还是候选行为；若不能重算，应限制跨版本比较。不应把旧规则下较低的起点与新规则下较高的终点直接相减，称作智能体能力提升。

评价标准也可能存在真实的价值或业务取舍。例如一种输出格式过去必须满足，后来变为可选。此时不是验证器程序出错，而是任务要求改变。新旧分数分别对各自要求有效，不能宣布某一时期的系统被事后证明更聪明。记录规则改变的原因和生效范围，有助于将能力变化与目标变化分开。

形式方法可以帮助验证特定性质，但仍依赖形式化规范与环境假设。如果性质只是“输出字段齐全”，证明它不等于证明答案正确；如果模型外部工具行为不在形式模型内，也不能把结论外推到整个服务。经验评价和形式检查可以组合，各自覆盖什么应清楚标注。本书目前没有为 E10 建立形式证明，不用自修改方法的名称暗示这项能力存在。

\subsection{恢复一个执行器，还是恢复整个演化过程}
\label{sec:ch12-restore}

第~\ref{ch14}章讨论运行系统的恢复；自修改还要恢复搜索状态。任务执行实现、修改程序、档案、父代关系、参与生成的持久记忆与实际剩余预算共同决定下一代从哪里开始。仅回退执行器、仍沿用新版搜索记录，是一个新的组合起点，应另记身份。

可以分两步核对：先检查任务入口、修改入口与依赖能否按快照加载，再确认下一次候选生成实际读取的档案与资源账。恢复不能退还已经消耗的请求；快照保存的是当时额度，实际可用额度还要扣除后续尝试。坏候选可以留在只读档案，但是否还能成为父代，应由明确的准入规则决定。外部副作用的处置沿用第~\ref{ch14}章的恢复边界。

\section{把元层问题变成可执行比较}
\label{sec:ch12-e10design}

目前 E10 只有教学级研究方向，本节将其拆成准备事项，不宣称协议已经冻结。核心问题可以先收窄为：在相同任务起点、开发反馈和资源范围内，允许修改候选生成程序，是否提高随后在留出任务家族上产生合格执行器的能力？这个问题比“系统是否学会无限持续改进”小得多，却具有可核对的对象。

静态参照冻结任务执行器和修改器，提供原系统的任务表现；执行器更新条件由固定修改器生成并采用新执行器；元层更新条件在相同执行器可改范围之外，再允许明确列出的修改程序进入候选，并在后续阶段实际使用。静态参照用于比较完整更新过程相对原系统的收益，后两个条件用于比较新增元层修改权限的作用。三者的任务起点、反馈范围和费用记录仍需明确，不能把静态参照没有发生的搜索费用补成与更新条件相同。

应先完成一个确定性的装置检查：用手写的合法候选确认发布后确实加载新版本，用不合法候选确认拒绝与恢复，用空改动确认不会被误计为机制创新。这能检查版本流转、评价隔离和预算，但不能进入模型效果表。若装置无法区分“文件写了”和“运行用了”，启动付费搜索只会放大记录问题。

随后冻结源领域和目标领域的用途。源领域提供允许的开发反馈，用来形成可迁出候选；目标开发范围允许修改器继续生成任务执行器；目标最终范围只用于已定选择后的评价。任务家族划分应防止模板变体跨界，并记录各领域是否共享工具、接口或基础知识。迁移有多远不能只由领域名称决定。

要评价修改器，就要保存失败候选。代码无法导入、权限越界、超时、预算中断、基本编辑能力丧失和最终任务表现不足，是不同失败原因。可报告有效候选比例与找到合格版本所用资源，但不应只对成功编译候选计算效果，然后忘记分母中所有未成功候选。发布门槛也需在看到结果前确定。

资源比较至少需要两个上限：产生候选的总模型与工具预算，以及评价这些候选的总预算。某个条件生成较复杂候选可能增加评价开销；若所有条件只限制候选数量，复杂路线就可能得到更多实际计算。可以选择成本配平或候选数配平，但应说明剩余差异，不把两者混为一谈。

本书现有小规模编程试跑不等于 E10。它们没有包含可变修改器的完整循环，也没有这里要求的源目标设计，因此不能拼接旧轨迹生成元层结果。E10若进入真实运行，需要独立的新协议、工件标识与相应资源授权；在这些条件未完成时，正文保留设计，不给预期收益或置信区间。

\subsection{一个候选从提出到采用的具体过程}
\label{sec:ch12-candidatewalk}

设想修改器发现执行器经常整文件覆盖，因而建议增加按唯一文本片段替换的编辑接口。这是教学案例，与来源论文中细粒度编辑的方向相近，但下面的验收步骤是本书提出的实现分析，不是对原作者某次运行日志的复述。

首先固定候选起点与修改范围。候选可以改编辑接口和调用文档，不能改评价题或日志采集包装层。修改器读取允许的失败案例后产生补丁；补丁保存来源、父代与生成费用。若它顺便修改预算常量，应视为范围外变化处理，而不能因为核心功能有用就忽略。

静态检查可以确认接口签名、路径处理和依赖能否加载，但它不能证明编辑语义正确。功能检查应包含唯一命中、没有命中、多次命中、空片段与换行差异等情况。是否允许多次命中必须由接口约定决定，不能让实现默默选择一个便利行为，再把所有调用者都当成同意了它。

下一步检查执行器是否真的使用新接口。工具定义可以已经更新，但提示仍只描述整文件写入；此时调用轨迹能揭示新功能未被采用。若采用后任务分未提高，也可能因为评价题不需要细粒度编辑，或额外选择开销抵消了节省。应把接口正确、实际使用与最终收益三个结论分开。

如果新接口参与自我修改，还要让候选执行器在受控副本中修改一个后续版本。任务修复通过不保证自修改路径正常：目录不同、模块加载方式不同或候选发布顺序，都可能触发额外问题。这个功能检查只是确认它还能参与下一代生成，不证明下一代一定更好。

通过开发检查后，仍需按预定规则选择是否发布。若旧任务出现回归，应回到具体失败而不是只看总均分；若新增工具节省调用但增加权限面，应按实际边界核对。发布记录最终绑定代码、文档、依赖、预算和评价结果，使后续能够复查这次接受的理由。

若候选被拒绝，它的生成和检查费用仍入账，失败原因进入只读档案。后续修改器可以在允许范围内学习这份失败记录，但不能把开发失败题当作新的最终封存题反复宣传。一次清楚的拒绝也是实验结果的一部分，它说明某个具体改动在当前条件下未达到接受要求。

\subsection{何时继续搜索，何时先修评价}
\label{sec:ch12-stopping}

发现候选开发分上升时，可以继续搜索；但若上升伴随日志缺失、评价版本不一致或数据划分失效，应先处理测量问题。继续运行会产生更多无法解释的分数，不能通过增加样本量修复系统性观测错误。保存当前工件和错误范围，比立即追求下一代更有价值。

连续多轮没有新最佳值，也不自动说明修改器失效。可能是预算太少、可改范围有限、评价噪声大、任务已接近上界，或搜索确实停滞。可以查看有效候选比例、费用与改动类型，决定是否继续既定预算、缩小问题或另立设计。若改变了搜索规则，应开启新阶段标记，不把干预前后当作同一个固定条件。

停止规则可以先以资源与质量约束表达：剩余预算不足以完成最小评价时停止生成；观测链不能核对时停止发布；候选失去基本功能时拒绝该候选；需要改变研究问题时记录新版本协议。它们不要求事先知道搜索的最佳终点，却能防止系统靠无限追加尝试来维持“最终总会提高”的叙述。

最后应交付一个具体版本，而不是一条抽象上升曲线。版本要有对应的可复查任务表现、资源记录、修改权限与已知失败范围；档案用于解释它如何产生。若没有候选达到预定要求，可以交付研究记录并保留原执行器。这样的结束同样完成了一次设计明确的试验，而不是把负结果隐藏在“继续演化中”。


\section{条件变化：元进化在什么时候退化}
\label{sec:ch12-conditions}

\begin{itemize}
  \item \textbf{基座模型已经很强时}，仍要实测额外脚手架的作用。
        SICA 的推理任务提升有限；独立 o3-mini 在 high 推理强度下达到 \cited{87\%} 与 \cited{79\%}，复合系统采用另一套调用安排，差值不能单独归给脚手架。
  \item \textbf{迁移到设计领域之外时}，应检查指令与工具假设。所引固定代码修改器在两个新领域的中位为\cited{0.0}，这是该实现的结果；手工适配与自动修改均可比较，不排除其他固定修改器。
  \item \textbf{去掉档案并改变父代选择时}，所引消融两个领域成绩分别为\cited{0.0}与\cited{0.116}。该结果支持具体组合的比较，不保证任何档案都能改善探索。
  \item \textbf{改动会碰到评测脚手架时}，应检查最终评价与候选写入的隔离，并保护观测链。若评价规则也要更新，需保留版本与共同外部参照，分开规则变化和候选行为变化。
\end{itemize}

% ---------- 8. 能说与不能说 ----------
\section{让自修改对应可检验的收益}
\label{sec:ch12-claims}

代码自修改把工具、执行过程与候选生成程序纳入持续改进的对象。SICA、DGM和HyperAgents分别提供具体的实现、搜索与迁移证据；它们的共同价值，是让“系统怎样改变自己”能够落实到版本、加载路径和后续行为。任务成绩、找到候选的效率与迁出后继续产生改进的能力，仍需分别测量。

HyperAgents的元层迁移结果为P7提供了正向外部证据，v2的组件对照进一步表明，在所研究的数学评分迁移配置中，主要增益随元Agent转移。这个结论对应具体实现、目标任务与搜索预算；本书E10尚未运行，仍需明确自身的任务族、迁移对象和资源范围。

\begin{principlebox}
元层可改范围扩大时，版本链和评价参照也要随之明确。用固定预算的候选生成与独立评价检验修改能力；默认将最终判据和观测记录置于候选写入之外。若另行研究评价器更新，应绑定规则版本、独立接受依据和可比较范围，而非把分数变化全部计为能力改善。
\end{principlebox}

\falsifier{P7需要先定义目标任务族、迁移对象、对照与资源范围。有效正向迁移可以反对对应范围的否定性主张；比较两种方法时，需要方法间差异及其不确定性，单组\texttt{imp@50}区间不能代替。E10未运行，相应条件仍是后续研究设计。}

\openquestion{当前没有本书自修改循环的实测。第13、14章的时机与恢复框架可作起点；候选包含修改器程序时，还需核对其依赖，恢复修改器、档案、谱系、持久状态与实际剩余预算，而非只退回任务执行代码。}

\section*{习题}
\begin{enumerate}[label=12-\arabic*.,ref=12-\arabic*]
  \item\label{ex:ch12-1} 表~\ref{tab:ch12-sica} 的两个自带基准完全定义在同一个可编辑代码库内。
        请列出还需核对的评价版本、读写权限与开发使用记录，再分别讨论：仅代码共处、实际修改了评价逻辑、测试题反复用于选择候选。这三种情况各支持什么结论？
  \item\label{ex:ch12-2} \S\ref{sec:ch12-archive} 的入档条件里有一条“必须保留代码库编辑能力”。
        若去掉这一条，请提出两种可能的档案变化，说明各自还依赖哪些选择规则，并设计一个能区分它们的观测量。
  \item\label{ex:ch12-3} 表~\ref{tab:ch12-imp}的\texttt{imp@50}由固定预算下选定后代的任务分数减去起点分数得到。
        请构造两组起点与后代分数，使任务分数整体提高而\texttt{imp@50}不变。
        这一观察能否识别修改器的贡献？还需哪种组件对照？
  \item\label{ex:ch12-4} \S\ref{sec:ch12-vex} 默认将最终判据与候选写入隔离。
        若某个系统必须让验证器也可更新（例如验证器随业务规则变化），
        请设计一种仍然能报告分数的安排，并说清它的代价。
  \item\label{ex:ch12-5} 为什么一个方法的 \texttt{imp@50} 大于零，不自动等于它比另一种方法更能迁移？请分别写出两种问题的比较对象。
\end{enumerate}

\section*{常见误解}
\begin{itemize}
  \item \textbf{“能改自己的系统一定比固定修改器的系统强”。}
        外部结果里 DGM-H 与人工按领域改造的差异并不显著（$p>0.05$）；
        多出来的一层提供了调整修改过程的可能性，是否改善任务表现与资源使用仍需对照。
  \item \textbf{“档案越大越好”。}
        入档条件与父代选择都在控制档案的形状；
        不检查基本功能，可能保留大量无法继续修改的节点；实际后果取决于选择规则，不能预言档案必然退化成某一种形状。
  \item \textbf{“自带基准的涨幅足以证明独立新任务能力”。}
        共处代码库提示权限风险，但修改执行代码不等于实际修改基准。应核对评价版本、读写权限和数据是否参与选择，再限定证据范围。
  \item \textbf{“元层再往上加一层就能解决固定修改器的问题”。}
        若新增的上层仍被固定，原问题会转移到新的边界；自指是原作者的一种应对设计，不据此断言所有分层方案均无效。
\end{itemize}

\section*{进一步阅读}
\begin{itemize}[itemsep=2pt,topsep=4pt,parsep=0pt]
  \item \cite{sica-2504}：Table 1 的逐迭代成绩、§3 的效用式与超时惩罚、
        §4 关于两个自带基准的说明（含作者自陈的基线数据质量问题），
        以及 §4.1 Figure 4 中独立推理模型与复合系统的不同调用条件。
  \item \cite{dgm-2505}：§3 的父代选择与入档条件，§4.1 与 §4.4 的成绩，
        §6 的单次运行约两周——这是“元进化的成本长什么样”最好的一处具体材料。
  \item \cite{hyperagents-2603}：Figure 1 的 DGM 与 DGM-H 对照（原论文图注为 DGM 与 DGM-H 的双栏对照，并非三层框架；本书三层框架见~图~\ref{fig:ch12-meta}）、Figure 2 的成绩与置信区间、
        Figure 3 与Appendix D.3—D.4的 \texttt{imp@50}及迁出选择规则；Appendix E.1展示任务侧改动。另读\cite{hyperagents-2603-v2} §5.2、附录D.5的Table 4，它在v1整包迁移基础上加入任务与元组件的交叉对照。
  \item \cite{memscope-2610}：§3 与附录A--B的局部提供者修补、无补丁决策，以及 Table 3 中旧成功、恢复与退化的配对计数。
  \item 附录~\ref{app:lab}：任务集划分与隔离边界在实现上落在哪些字段。
  \item 附录~\ref{app:cost}：区分整次搜索、单次迭代与最后版本的执行成本；原文约两周对应完整运行。
\end{itemize}
